\section{Discussion}\label{sec:discussion}

Robotic active matter has developed into an experimental
field that connects programmable interactions, collective
motion, and nonequilibrium statistical physics. Its central
achievement is not the reproduction of patterns that resemble
those of biological or synthetic active systems. Rather, it
is the ability to modify microscopic interactions and observe
their collective consequences in physical hardware.

Yet, the evidence reveals an important asymmetry.
Experimental demonstrations of collective behavior are
widespread, while independently replicated quantitative
tests of general physical theories remain less common.
This difference is central to our interpretation of
the five review objectives.

For objective O1, we find that robotic active matter
encompasses several distinct experimental regimes:
programmed mobile swarms, mechanically driven particles,
microscopic robotic agents, feedback-controlled colloids,
modular assemblies, and active mechanical structures.
These platforms make different physical quantities
accessible. Their diversity is an experimental resource,
but it does not imply that their collective dynamics
follow the same theoretical description.

For O2, we find broad experimental evidence for
flocking, clustering, transport, collective assembly,
and several forms of rotational and synchronized
motion. More specialized phenomena, including
active elasticity and odd mechanical response,
have also been investigated. However, a common
phenomenological label can conceal fundamentally
different mechanisms.

For O3, quantitative theory--experiment comparisons
are present across several frameworks, but
agent-based simulation remains the dominant
comparator. The distinction between reproducing
an implemented controller and predicting a
collective physical observable is therefore
essential.

For O4, sensing, communication, and feedback
make collective interactions experimentally
adjustable. However, the controller is itself
part of the physical dynamics and can impose
the behavior that a theoretical model
purports to explain.

Finally, O5 reveals substantial weaknesses
in reproducibility documentation, finite-size
characterization, and the reporting of
theoretical tests. These weaknesses do not
invalidate the reported observations.
They limit the strength and generality
of the conclusions that can be drawn
from them.

In the following, we assess these limitations
systematically and distinguish conclusions
that are strongly supported from those
that remain conditional.

\subsection{Risk of bias across the evidence base}
\label{sec:discussion-risk}

A central difficulty in reviewing robotic active matter
is the absence of a validated quality instrument
covering its heterogeneous experimental designs.
A microscopic colloidal experiment, an aerial
robot swarm, and an active mechanical metamaterial
have different measurement and control requirements.
Nevertheless, they share several conditions
necessary for reliable inference.

We therefore used six prespecified domains:
experimental replication (D1), agent
heterogeneity (D2), finite-size and boundary
effects (D3), measurement validity (D4),
data and code availability (D5), and
quantitative theory comparison (D6).

These assessments describe limitations in
the evidence supporting a scientific claim.
They are not measurements of experimental
correctness. In particular, an incomplete
report does not establish that the
corresponding experimental practice
was absent.

The evidence base contains 491 records.
The domain totals below refer to the
corrected October 17, 2026 extraction.
One record lacks all six domain ratings.

\paragraph*{Replication and independent runs.}

For D1, 126 studies are rated low risk,
348 have some concerns, and 17 are
high risk. Thus, among the 491 records
with a D1 rating, only 25.7\%
receive a low-risk assessment.

The most frequent limitation is incomplete
documentation of independent experimental
repetitions. A numerical run count is
recorded for 196 of the
491 studies (39.9\%).

This finding matters because active-matter
experiments often involve stochastic
trajectories, fluctuating configurations,
and finite populations. A collective state
observed in one realization may be sensitive
to initial conditions or environmental
fluctuations.

For example, experiments on collective
aggregation, active mechanical response,
and controlled colloidal organization
can provide strong physical measurements,
but reproducibility depends on how
consistently the reported observables
can be obtained
\cite{Garnier2005aggregationbehaviou,Buerle2020formationofstable,Heuthe2025tunablecolloidalsw,Arbel2025amechanicalroutef,Boudet2025nonmarkoviandynami,Chen2024persistentandrespo,Wang2025hyperuniformnetwork}.

A reported independent-run count is useful,
but it does not by itself establish
independence of every measurement.
Repeated trajectory samples from one
experimental run may be correlated.
Conversely, repeated experiments may
share calibration errors or
systematic hardware differences.

The important quantities are therefore
the number of independent realizations,
the variation between them, and the
uncertainty of the collective observable.
The extracted dataset does not provide
a standardized measure of all three.

Hence, the D1 results support a
moderate level of confidence in
experimental reproducibility across
the evidence base, with substantial
uncertainty for studies lacking
explicit repetition information.

\paragraph*{Agent heterogeneity.}

For D2, 389 studies are rated low risk,
98 have some concerns, and four
are high risk.

This is the most favorable distribution
among the six domains. It suggests that
agent properties are comparatively
well characterized in the assessed
experimental reports.

Yet, the meaning of heterogeneity
depends on platform. For wheeled
robots, relevant variations include
motor response, sensing range,
and turning accuracy. For vibration-driven
robots, they include propulsion
speed, body geometry, and
substrate-dependent friction.
For microscopic systems, particle
size, magnetic response, and
actuation efficiency may matter.

These differences can affect collective
behavior even when the nominal
interaction rule is identical.

Studies involving mechanically driven
particles, microscopic assemblies,
and programmed swarms illustrate
the range of relevant
heterogeneity mechanisms
\cite{Welch2018aggregatesuppressio,Deblais2019howtoforceanarmy,Xie2020bioinspiredsoftmic,Pauer2021programmabledesign,Li2019particleroboticsba,Vsrhelyi2018optimizedflockingo,Zion2023morphologicalcomput}.

The favorable D2 assessment should
therefore not be interpreted as
evidence that heterogeneity is
physically irrelevant. Rather,
it indicates that the extracted
reports generally provide
better information in this
domain than in several others.

A future quantitative synthesis
would benefit from recording
distributions of agent properties,
rather than only representative
values. This would permit direct
tests of whether heterogeneity
changes collective order,
transport, or stability.

\paragraph*{Finite-size and boundary effects.}

D3 is one of the most problematic domains.
Only 39 studies receive low-risk ratings,
445 have some concerns, and seven
are high risk.

Thus, only 7.9\% of records with a D3
rating are assessed as low risk.
This represents a major limitation
for conclusions about collective
phases and continuum behavior.

Many robotic experiments take place
in finite arenas containing
relatively small populations.
Under these conditions, boundaries
may influence a substantial
fraction of the agents.
The observed state can
therefore differ from
the behavior predicted
for an extended system.

Experiments investigating confinement,
spatial correlations, and collective
mechanical response demonstrate
why these effects matter
\cite{Deblais2018boundariescontrolc,Boudet2025effectsofboundarie,Casiulis2025ageometricconditio,Wang2025hyperuniformnetwork,Zhang2026emergentnonhermiti,Coulais2025activemetamaterials,Wagner2026geometryinducedtra}.

However, the extraction does not establish
that every such study performed a
complete finite-size analysis.
A measurement of boundary accumulation,
for example, identifies an important
physical effect, but does not
necessarily determine its influence
on a bulk phase transition.

The consequence is particularly
important for claims involving
long-range order, phase separation,
and continuum constitutive laws.

A finite group may exhibit persistent
polarization without establishing
thermodynamic long-range order.
A dense aggregate may resemble
phase coexistence without
demonstrating a bulk instability.
An active mechanical network
may display localized response
without establishing a
size-independent material property.

Hence, D3 limits the certainty
of broad theoretical conclusions
more strongly than it limits
the validity of individual
experimental observations.

\paragraph*{Measurement validity.}

For D4, 380 studies are rated low risk,
106 have some concerns, and five
are high risk.

This relatively favorable distribution
indicates that the extracted
measurement procedures are generally
more thoroughly documented than
finite-size analysis or
quantitative theory validation.

Robotic systems offer substantial
measurement advantages.
Macroscopic agents can often
be tracked directly.
Microscopic systems permit
optical observation of
particle trajectories.
Active structures can be
characterized through deformation,
spectral, or mechanical response.

Nevertheless, measurement validity
is observable-dependent.
A tracking procedure adequate
for center-of-mass displacement
may not resolve short-range
velocity correlations.
A method suitable for
visualizing assembly geometry
may not accurately measure
forces or fluctuations.

Relevant studies demonstrate
the range of experimentally
accessible observables
\cite{Tarcai2011patternstransition,Buerle2020formationofstable,Chvykov2020lowrattlingapred,Brandenbourger2021limitcyclesturnac,Barone2024experimentalandnum,Hanson2025electrokineticpropu,Engbring2023anonlinearfluctuat}.

We therefore interpret the D4 ratings
as evidence about the documented
validity of measurements, not as
a universal statement that
all reported observables
are directly comparable.

A stronger future extraction
would record measurement uncertainty,
spatial and temporal resolution,
tracking failure rates, and
validation against independent
measurement methods.

\paragraph*{Data and code availability.}

D5 presents the clearest reporting deficit.
Only 38 studies are rated low risk,
57 have some concerns, and 396
are high risk.

Explicit code availability is identified
in 71 of the 491 studies
(14.5\%). Data availability is
identified in 54 (11.0\%).

These fields were coded positively
only when a repository or archive
was named. Consequently, a negative
value does not necessarily imply
that the authors would refuse
to share their materials.
It indicates that availability
was not established under
the extraction rule.

The scientific consequences differ
between experimental platforms.

For programmed swarms, controller
code defines the effective
interaction law.
For feedback-controlled colloids,
the measurement and feedback
implementation determine
state-dependent interactions.
For vibration-driven robots,
trajectory data and calibration
measurements constrain
stochastic physical models.
For active solids, response
measurements and analysis
procedures determine the
estimated mechanical coefficients.

Without accessible materials,
independent researchers may be
unable to reproduce the
theory--experiment comparison
even when the published
qualitative result is clear.

Hence, the low availability
rates are a substantial
limitation of the
current evidence base.
They also identify a
practical opportunity:
standardized sharing of
trajectories, calibration,
control implementations,
and analysis scripts
could improve reproducibility
without requiring
new experimental hardware.

\paragraph*{Quantitative theory comparison.}

For D6, 83 studies are rated low risk,
146 have some concerns, and 262
are high risk.

Only 16.9\% of the 491 rated records
therefore receive a low-risk
assessment for quantitative
theory comparison.

The extracted agreement field provides
a related but different measure.
Of the 491 studies,
126 report quantitative agreement,
238 qualitative agreement,
124 no comparison, and three
have missing values.

The difference between these distributions
is important. A study may report
quantitative agreement while
the quality assessment identifies
limitations in the comparison.
For example, parameters may have
been fitted to the same observations
used for validation, or the
comparison may lack sufficient
uncertainty estimates.

Conversely, a qualitative comparison
can be scientifically valuable
when it identifies a new mechanism
or motivates a discriminating
experiment. It simply supports
a different level of inference.

The strongest theoretical evidence
includes experiments that connect
measured microscopic interactions
to collective observables
\cite{Buerle2020formationofstable,Xiao2023perceptionofmotion,Arbel2025amechanicalroutef,Heuthe2025tunablecolloidalsw,Boudet2025nonmarkoviandynami,Chen2024persistentandrespo,Chvykov2020lowrattlingapred,Wang2025hyperuniformnetwork,Zion2023morphologicalcomput,Binysh2026moreislessinunpe}.

Yet, even within this literature,
the extraction does not systematically
distinguish independent predictions
from fitted descriptions.

The central conclusion of D6 is
therefore not that existing
models generally fail.
Rather, the literature often
does not provide sufficiently
discriminating tests to
establish how far the
models can be trusted.

\paragraph*{Joint limitations.}

The six risk-of-bias domains should
not be interpreted independently.
A quantitative comparison becomes
less convincing when repetitions
are poorly documented.
Similarly, agreement with a
continuum model may be limited
by finite size even when
the experimental observable
is measured accurately.

Only 30 studies in the corrected
extraction receive low-risk
ratings in both D1 and D6.
This is 6.1\% of the
491-study evidence base.

These studies span several
experimental approaches
\cite{Garnier2005aggregationbehaviou,Buerle2020formationofstable,Yang2022emergentmicroroboti,Xiao2023perceptionofmotion,Arbel2025amechanicalroutef,Heuthe2025tunablecolloidalsw,Boudet2025nonmarkoviandynami,Casiulis2025ageometricconditio,Chen2024persistentandrespo,Chvykov2020lowrattlingapred,Savoie2019arobotmadeofrobo,Wang2025hyperuniformnetwork,Zion2023morphologicalcomput,Li2020programmingactiveg}.

This subset provides a useful starting
point for sensitivity analysis.
However, low risk in two domains
does not guarantee low risk
in the remaining four.
Nor does it establish
independence between
research groups.

We therefore avoid classifying
the entire literature as
either reliable or unreliable.
The appropriate judgment
depends on the specific
physical claim and the
evidence required to
support it.

\subsection{Reporting bias and publication status}
\label{sec:discussion-reporting}

Publication status provides one possible indicator
of reporting differences, but it is not a
direct measure of scientific quality.
Preprints may contain complete experimental
methods and quantitative comparisons.
Peer-reviewed articles may also omit
important details.

The review protocol specifies a comparison
between preprint-only and peer-reviewed
findings. However, the supplied extraction
contains a venue field rather than a
reconciled publication-status variable.
It does not establish whether a record
listed as an arXiv or bioRxiv publication
has a corresponding peer-reviewed version
outside the included table.

We can therefore report a descriptive
comparison of venue labels, but not
a definitive preprint-only versus
peer-reviewed analysis.

In the original extraction,
47 records have venues explicitly
containing arXiv or bioRxiv.
The remaining 445 have other
venue labels or missing
venue information.

Among the 47 preprint-labelled records,
18 report quantitative theoretical
agreement, 21 qualitative agreement,
and eight no comparison.
Their corresponding proportions
are 38.3\%, 44.7\%,
and 17.0\%.

Among the remaining 445 records,
109 report quantitative agreement,
217 qualitative agreement,
116 no comparison,
and three have missing values.
Quantitative agreement is
recorded for 24.5\%
of this group.

Thus, the extracted preprint-labelled
records contain a higher proportion
of quantitative comparisons.
Yet, this difference cannot
be interpreted as a
publication-bias effect.

The groups differ in disciplinary
composition, publication year,
and potentially the nature
of the experimental platforms.
Furthermore, the comparison
does not reliably separate
preprint-only studies from
preprints associated with
peer-reviewed publications.

For D6, ten preprint-labelled
studies are rated low risk,
20 have some concerns,
and 17 are high risk.
Among the other records,
73 are rated low risk,
126 have some concerns,
245 are high risk,
and one has no rating.

For replication, nine
preprint-labelled studies
are low risk,
37 have some concerns,
and one is high risk.
Among the remaining records,
117 are low risk,
311 have some concerns,
16 are high risk,
and one lacks a rating.

These descriptive differences
show that the preprint-labelled
subset is not uniformly weaker
than the remainder of the
literature under the
extracted quality criteria.

Nevertheless, the comparison
cannot establish whether
positive collective phenomena
are preferentially published.
The dataset contains included
studies only. It does not
provide a representative
sample of unsuccessful
experiments, unpublished
null results, or
abandoned theoretical
comparisons.

Publication bias may therefore
operate before a study enters
the searchable literature.
A collection of reported
collective states can appear
more robust than the
underlying experimental
research process if
negative findings are
less frequently disseminated.

A further problem is the
selective reporting of
observables. One experiment
may measure several order
parameters but publish
only those showing
a clear effect.
The extraction cannot
detect such selection
without protocols,
complete datasets,
or additional
experimental records.

Hence, we cannot assign
a reliable numerical
publication-bias correction.
The evidence supports
a qualitative concern
about selective reporting,
not a quantitative
estimate of its magnitude.

A stronger future analysis
would require reconciliation
of preprints with their
published counterparts
and explicit extraction
of publication history.
It would also benefit
from complete experimental
datasets containing
unsuccessful or
null-condition runs.

\subsection{Certainty of the main conclusions}
\label{sec:discussion-certainty}

We now translate the evidence into certainty
judgments for the principal scientific conclusions.
Following the prespecified framework, we use
high, moderate, and low certainty.
These ratings consider experimental breadth,
independent support, replication,
and the D1 and D6 assessments.

The ratings apply to particular statements.
They should not be interpreted as a
single quality score for an entire
platform or phenomenon family.

\paragraph*{Robots can realize collective active behavior: high certainty.}

The existence of physical robotic collectives
exhibiting coordinated motion, spatial
organization, transport, and collective
reconfiguration is supported by
multiple experimental platforms.

Evidence includes programmed mobile robots,
mechanically driven particles,
microscopic assemblies, and
modular systems
\cite{Rubenstein2014programmableselfas,Vsrhelyi2018optimizedflockingo,Welch2018aggregatesuppressio,Li2019particleroboticsba,Yigit2019programmablecollect,Buerle2020formationofstable,Savoie2019arobotmadeofrobo,Xie2020bioinspiredsoftmic,Yang2022emergentmicroroboti,Xi2024emergentintelligenc,Wang2025hyperuniformnetwork}.

These observations do not depend
on one theoretical interpretation.
The systems physically exhibit
the reported collective responses.

We therefore assign high certainty
to the broad conclusion that
robotic platforms can realize
measurable collective active
behavior.

This rating does not extend
automatically to the mechanisms
proposed to explain each state.

\paragraph*{Collective mechanisms generalize across platforms:
moderate certainty.}

Several collective phenomena occur in
systems with substantially different
propulsion and interaction mechanisms.
This supports the view that some
principles of active collective
behavior extend beyond
individual hardware designs.

Relevant examples include
collective ordering,
mechanically mediated transport,
controlled clustering,
and collective rotation
\cite{Tarcai2011patternstransition,Scholz2018rotatingrobotsmove,Buerle2020formationofstable,Chvykov2020lowrattlingapred,Fruchart2022nonreciprocalphase,Chen2024persistentandrespo,Luo2025flockingphasesepar,Casiulis2025ageometricconditio}.

However, similar phenomena
can arise through different
microscopic interactions.
Moreover, normalized
information about independent
experimental groups is
not available.

We therefore assign
moderate certainty to
the general conclusion
that related collective
mechanisms can be
realized across
multiple platforms.

We do not infer
that all implementations
belong to the same
universality class.

\paragraph*{Specific universal phase behavior:
low certainty.}

Claims about universal flocking transitions,
motility-induced phase separation,
or thermodynamic scaling require
measurements beyond the existence
of an ordered or clustered state.

The limited number of
low-risk D3 ratings
and the prevalence
of qualitative comparisons
restrict confidence
in such conclusions.

Although individual studies
provide quantitative evidence
\cite{Tarcai2011patternstransition,Buerle2020formationofstable,Xiao2023perceptionofmotion,Barone2024experimentalandnum,Wang2025hyperuniformnetwork,Verdoucq2025flockingphasetrans},
the supplied extraction
does not establish
consistent finite-size
scaling or independent
validation across
the entire literature.

We therefore assign
low certainty to
broad claims that
specific universal
phase behavior
has been established
across robotic
active-matter
platforms.

This judgment does
not exclude stronger
conclusions for
particular models
and experimental
conditions.

\paragraph*{Quantitative theory--experiment agreement:
moderate certainty.}

Quantitative comparisons are documented
for Vicsek-type models,
active Brownian particles,
kinetic and hydrodynamic theories,
active elasticity,
and other statistical-mechanical
descriptions.

Several studies demonstrate
the feasibility of connecting
microscopic dynamics to
measured collective observables
\cite{Buerle2020formationofstable,Chvykov2020lowrattlingapred,Brandenbourger2021limitcyclesturnac,Xiao2023perceptionofmotion,Barone2024experimentalandnum,Heuthe2025tunablecolloidalsw,Boudet2025nonmarkoviandynami,Hanson2025electrokineticpropu,Zion2023morphologicalcomput}.

The evidence supports
moderate certainty
that quantitative
active-matter models
can describe selected
robotic experiments
successfully.

Yet, the extracted
agreement labels do
not systematically
distinguish fitted
from independently
predicted observables.
Therefore, the broader
claim that these
models possess
general predictive
validity across
platforms receives
low certainty.

\paragraph*{Programmability provides causal experimental control:
high certainty.}

Robotic platforms permit
controlled modification
of interaction rules,
feedback architecture,
sensing, and actuation.

This capability is
demonstrated across
macroscopic swarms,
feedback-controlled
colloids, and active
mechanical structures
\cite{Ferrante2013aselfadaptivecomm,Li2019particleroboticsba,Buerle2020formationofstable,Chen2022collectiveresponse,Fruchart2022nonreciprocalphase,Heuthe2025tunablecolloidalsw,Brandenbourger2021limitcyclesturnac,Coulais2025activemetamaterials}.

We assign high certainty
to the conclusion
that programmability
makes interaction
parameters experimentally
accessible.

However, the stronger
claim that programmability
has already enabled
decisive tests of
general active-matter
theories receives
only moderate certainty.

The distinction follows
from the limited
number of low-risk
quantitative comparisons
and the possibility
that controllers
directly impose
the observed behavior.

\paragraph*{Experimental reproducibility is adequately documented:
low certainty.}

The extracted reporting indicators
show substantial missingness
in repetition counts,
data availability,
and code availability.

Only a minority of
studies receive
low-risk ratings
in both replication
and quantitative
theory comparison.

We therefore have
low certainty that
the currently
available literature
provides sufficient
documentation for
broad independent
reproduction of
its quantitative
claims.

This is not a
claim that the
experiments themselves
are generally
irreproducible.
It is a conclusion
about what the
available reports
permit us to verify.


\subsection{Limitations of the evidence}
\label{sec:discussion-evidence-limitations}

The certainty judgments above must be interpreted in light
of several limitations that extend beyond the individual
risk-of-bias domains. These limitations concern the
composition of the evidence base, the comparability
of experimental platforms, the distinction between
collective phenomena and physical mechanisms, and
the way experimental findings enter the literature.

The corrected October 17, 2026 snapshot comprises
491 included studies. Its breadth is a strength,
but it is not a representative sample of every
experimental attempt to realize collective
active behavior. Studies enter the evidence base
through publication, indexing, retrieval,
eligibility assessment, and extraction.
Each step can affect which phenomena
appear most strongly supported.

\paragraph*{Heterogeneity of physical implementations.}

The included platforms differ substantially in
propulsion, sensing, interactions, agent size,
collective size, and control architecture.
These differences are scientifically
interesting, but they limit direct pooling
of experimental results.

For example, collective ordering in a
communicating mobile swarm
\cite{Turgut2008selforganizedflocka,Vsrhelyi2018optimizedflockingo,Zhou2024anticipationpromote}
does not necessarily arise from the same
effective interactions as ordering
in mechanically driven robots
\cite{Welch2019activereorientation,Baconnier2022selectiveandcollec,Luo2025flockingphasesepar}
or feedback-controlled colloids
\cite{Buerle2020formationofstable,Chen2024synchronousandfull,Chen2024persistentandrespo}.

The resulting order parameters may be
similar, but the assumptions governing
interaction range, conservation laws,
noise, and information transfer differ.

This creates a distinction between
phenomenological replication and
mechanistic replication.
Independent experiments may reproduce
a visually similar state without
independently confirming the same
theoretical explanation.

Hence, our synthesis emphasizes
observables and interaction
mechanisms rather than treating
every report of a phenomenon
as equivalent evidence.

\paragraph*{Small systems and finite experimental arenas.}

A further limitation is experimental scale.
Many robotic studies use small collectives
because individual agents are expensive,
require tracking or calibration,
or operate within restricted arenas.

Small systems are valuable for investigating
pairwise coupling, local response,
and controlled interaction rules.
They are less decisive for conclusions
about collective phases that depend
on system-size scaling.

For example, stable polarization
in a small flock does not
establish long-range order
in an extended active system.
A dense aggregate in a confined
arena does not by itself
demonstrate macroscopic
phase coexistence.

Similarly, mechanically coupled
robotic assemblies can display
nonreciprocal response or
localized motion without
establishing a size-independent
constitutive law
\cite{Brandenbourger2019nonreciprocalrobot,Brandenbourger2021limitcyclesturnac,Binysh2026moreislessinunpe,Coulais2025activemetamaterials}.

The unfavorable D3 assessment
across the evidence base
therefore restricts generalization
from finite experimental
realizations to universal
active-matter behavior.

An important consequence is
that increasing the number
of studies does not
automatically resolve
this limitation.
Many reports using
similarly small systems
may leave the same
scaling question unanswered.

\paragraph*{Outcome definitions and theoretical specificity.}

The phenomenon taxonomy is intentionally broad.
It accommodates the terminology used by
different experimental communities.
However, it also combines outcomes
that are not theoretically equivalent.

The category of clustering or MIPS
includes experiments with programmed
attraction, field-mediated assembly,
contact-driven aggregation,
and other collective structures.
Only some can directly test
motility-induced phase separation.

Likewise, the categories of
chirality or rotation,
swarming or vortices,
and self-assembly or shape formation
cover several distinct mechanisms.

Experiments on programmed assembly
\cite{Rubenstein2014programmableselfas,Li2019particleroboticsba,Yi2025reconfigurablerobot}
and microscopic field-driven
collectives
\cite{Pauer2021programmabledesign,Sun2024multiplemagnetoopt,Cheng2025selfassemblyandco}
may therefore share a
phenomenon label while
testing different
physical questions.

This limitation is particularly
important for quantitative synthesis.
A common label does not justify
combining order parameters,
effect sizes, or agreement
statistics without additional
harmonization.

We consequently avoided a
pooled effect estimate.
The absence of such an estimate
is not merely a consequence
of missing numerical data.
It reflects genuine differences
in mechanisms, experimental
conditions, and estimands.

\paragraph*{Independence of experiments and research groups.}

Publication counts are not equivalent to
independent experimental confirmations.
Several reports may use related hardware,
similar controllers, or data generated
within the same research program.

The extraction provides citation keys,
publication metadata, platform classes,
and experimental descriptors.
It does not contain a fully reconciled
mapping of shared apparatus,
experimental lineages,
or independent research groups.

We therefore cannot reliably count
independent replications of a
particular theoretical claim.

This limitation affects
certainty judgments even
when a phenomenon appears
in many publications.
Repeated demonstrations
within one experimental
tradition provide useful
evidence for feasibility,
but they do not replace
independent tests using
different apparatus
and measurement pipelines.

Cross-platform evidence partially
addresses this concern,
yet platforms are not
equivalent to research groups.
Nor does hardware diversity
guarantee that the same
theoretical prediction
was tested.

\paragraph*{Null results and selective outcome reporting.}

The available literature is likely
to contain a disproportionate
number of successful demonstrations
of collective behavior.
This concern cannot be
quantified from the
included studies alone.

The evidence base does not
provide a registry of
attempted experiments,
unsuccessful implementations,
or unreported parameter sweeps.

Consequently, we cannot
estimate how often
a proposed interaction
fails to generate
the expected state.

Selective reporting may also
occur within a study.
An experiment may measure
several collective observables
while emphasizing only
those that exhibit
the clearest response.

This possibility is
especially relevant when
a model has several
adjustable parameters
or when a controller
is optimized to produce
a particular pattern.

The absence of complete
datasets and analysis
implementations makes
these reporting effects
difficult to assess
retrospectively.

We therefore avoid
interpreting the
frequency of positive
agreement labels
as an unbiased
estimate of
the probability
that a theoretical
model is correct.

\paragraph*{Concentration of theoretical evidence.}

The theoretical comparisons are not
distributed evenly across
frameworks or platforms.

Agent-based simulation is
widely used, whereas
some analytically
specific descriptions,
including run-and-tumble
models and active elasticity,
are represented by
much smaller literatures.

This unevenness limits
comparisons between
theories.
A framework with
many implementation-based
comparisons is not
necessarily better supported
than one with fewer
but more discriminating
quantitative tests.

The relevant distinction
concerns independent parameter
constraints, measurement
uncertainty, and the
ability of a model
to predict results
outside its fitting
conditions.

The available extraction
records agreement type,
but it does not
systematically separate
fitted and predicted
quantities for every
included study.

This prevents us
from assigning
a numerical
predictive-validity
score across
theoretical frameworks.

\subsection{Limitations of this review}
\label{sec:discussion-review-limitations}

The review also has limitations arising from its
search strategy, eligibility rules, automated
assistance, extraction process, and treatment
of incomplete reports. These limitations
must be distinguished from shortcomings
in the underlying experiments.

We used a prespecified protocol,
PRISMA-based selection accounting,
a controlled extraction vocabulary,
and documented validation and
adjudication. These procedures
improve transparency.
They do not eliminate uncertainty.

\paragraph*{Database coverage and retrieval.}

The search identified 15\,034 database records.
After removing 6\,332 duplicates,
8\,702 records were screened.
The selection process ultimately
produced the corrected evidence
base of 491 studies after
one included duplicate
was merged.

Database searching cannot guarantee
complete coverage of robotic
active matter.
The field crosses
physics, robotics,
materials science,
and applied engineering.
Relevant experiments
may use different
terminology depending
on their disciplinary context.

For example, an experiment
investigating active
mechanical response
may not use the
phrase ``active matter.''
A distributed robotic
system may report
collective fluctuations
while being described
primarily as
multi-agent control.

Search strings can be
expanded to improve
coverage, but broadening
them also increases
the number of irrelevant
records requiring
assessment.

The protocol's requirement
for physical interacting
artificial agents and
a collective or
interaction-level
observable also defines
the scope of this review.

Simulation-only studies
without the specified
hardware connection
were excluded.
So were studies
concerned solely
with task performance
when no qualifying
collective observable
was established.

These exclusions were
necessary to answer
the review questions.
However, they mean
that the review
does not represent
the entire theoretical
or engineering
literature on
collective robotics.

We do not infer
that excluded
studies are
scientifically
unimportant.
They fall outside
the defined
evidence base.

\paragraph*{Incomplete retrieval and abstract-only assessment.}

The selection flow distinguishes
reports assessed on full text
from reports assessed only
on abstracts because full
text could not be obtained.

In the reconciled
pre-merge selection flow,
979 reports were
assessed on full text:
462 were included,
513 excluded,
and four awaited
a decision.

Another 179 reports
were assessed on
their abstracts
because full text
was unavailable:
30 were included
and 149 excluded.

The resulting
662 eligibility
exclusions comprise
513 full-text
exclusions and
149 abstract-only
exclusions.

The 491 included records
therefore consist
of 461 full-text-based
and 30 abstract-only
inclusions.

After the subsequent
duplicate merge,
the corrected
included-study count
is 491.
The identity and
assessment basis
of the merged
duplicate have not
been specified,
so we do not
assign the resulting
reduction to either
assessment stratum.

The four reports
awaiting a decision
remain outside
the included
evidence base.

Abstract-only assessment
creates several
distinct uncertainties.
An abstract may
establish that an
experiment involved
physical agents
and collective behavior,
but it may not
describe experimental
repetitions,
measurement validation,
control implementation,
or the extent
of theoretical
agreement.

The extracted fields
for such studies
are therefore
particularly susceptible
to missingness
and incomplete
characterization.

An abstract may also
use broad terminology
that cannot be
fully resolved
without the
methods and results.

We retained these
studies under
the prespecified
eligibility process,
while avoiding
unsupported detailed
claims where the
available evidence
is insufficient.

Nevertheless,
their presence
limits the precision
of some descriptive
statistics and
certainty assessments.

A useful future
sensitivity analysis
should compare
conclusions with
and without
abstract-only records.
Such an analysis
would not recover
the missing details,
but it would
establish whether
the broad synthesis
depends materially
on these inclusions.

\paragraph*{LLM-assisted screening and extraction.}

The review used large language models
to assist screening, extraction,
and parts of the adjudication
workflow. This approach enabled
a substantial literature
to be processed using
a structured protocol,
but it introduces
additional sources
of error.

An automated screener
may misinterpret
whether a study
contains physical
hardware, whether
multiple agents
interact, or whether
the reported
observable is genuinely
collective.

Similarly, an extraction
system may assign
a platform or
phenomenon label
based on terminology
that does not
fully describe
the experimental mechanism.

These errors can
be systematic.
For example,
a model may
favor familiar
platform categories
or interpret
collective task
performance as
evidence of
collective active
dynamics.

The use of
structured vocabularies
reduces some ambiguity,
but it can also
produce forced
classifications
when the available
information is
incomplete.

We therefore
regard automated
extraction as
an evidence-processing
procedure requiring
validation, rather
than as an independent
source of
scientific authority.

\paragraph*{Blind validation versus final-label validation.}

The validation exercise
provides important evidence
about the performance
of the screening procedure,
but it must be
interpreted according
to which labels
were used as
the reference.

The first author
labels were assigned
blindly in a
200-record sample.
Against these
labels, the
review panel's
measured performance
was substantially
less favorable
than against
the final
adjudicated labels.

The initial
blind-label comparison
for the panel
gave sensitivity
of 0.290,
specificity
of 0.917,
and Cohen's
$\kappa=0.232$.

The final-label
validation produced
sensitivity of 1.000,
specificity
of 0.983,
and $\kappa=0.922$.

These are not
interchangeable
performance estimates.

The first comparison
reflects agreement
with the initial
blind judgments.
The second reflects
agreement with
the labels established
after adjudication
and review.

The improvement
is relevant
to the quality
of the final
classification process.
However, the
final-label comparison
does not constitute
a fresh independent
blind validation
of the same
system.

In particular,
the reference
labels changed
during the
adjudication process.
The high final
agreement therefore
cannot be interpreted
as if it arose
from an untouched
external gold
standard.

We use the
two comparisons
for different
purposes.
The initial
blind exercise
identifies disagreement
and ambiguity.
The final-label
analysis characterizes
the agreement
of the resulting
screening output
with the adjudicated
reference.

Both are informative,
but neither
eliminates the
possibility of
shared errors
between automated
systems and
human adjudication.

\paragraph*{The Opus judge and adjudication dependence.}

The adjudication process also used
Claude Opus as a judge
for records with
conflicting screening
decisions.
This role was
separate from
the model's blind
screening evaluation
on the validation sample.

For 423 full-text
screening disagreements,
the judge agreed
with the panel
in 362 cases,
reversed the
panel's decision
in 59,
and left two
undecided.

These results
show that
adjudication had
a material effect
on the final
evidence base.

They should not
be interpreted
as an independent
accuracy estimate
for the judge.
The absence of
a universally
available external
reference prevents
such a conclusion.

The judge also
consulted full text
directly for
a subset of
disagreements.
In other cases,
the decision
depended on
the evidence
available within
the adjudication
workflow.

This creates
the possibility
that different
disagreement types
received different
amounts of
source verification.

The protocol
and adjudication
records document
these decisions,
which supports
auditability.
Nevertheless,
a model serving
as judge may
share systematic
interpretive tendencies
with the models
whose output
it adjudicates.

The validation
sample provides
an additional
perspective.
As a blind
screener against
the final
author labels,
Opus had
sensitivity
of 0.550,
specificity
of 1.000,
and Cohen's
$\kappa=0.687$.

These values
must not be
combined with
the adjudication
agreement rate
as though
they measured
the same task.
Blind screening
and adjudication
of known disagreements
have different
information conditions.

A further limitation
concerns extraction
adjudication.
The extraction
workflow identified
167 conflicting
fields, with
18 ultimately
set to null.
These unresolved
or insufficiently
supported values
contribute to
missingness in
the synthesis.

The null values
are preferable
to invented
precision.
However, they
also restrict
the conclusions
that can
be drawn
from certain
records.

\paragraph*{The residual platform category.}

The platform taxonomy contains
a substantial residual
``other'' class:
88 of the corrected
491 studies,
or approximately
17.9\%.

This category includes
physically different
systems that do
not fit the
prespecified
platform classes.

It therefore
cannot be treated
as a coherent
experimental mechanism.

The residual category
has two consequences.
First, its size
limits the
resolution of
platform-based
comparisons.
Second, it
indicates that
the taxonomy
may need
further refinement
as new experimental
implementations
enter the field.

Nevertheless,
retaining an
``other'' category
is preferable
to forcing
studies into
classes that
do not accurately
describe their
hardware.

The same principle
applies to
the four records
without assigned
platform classes.
Missingness
should remain
visible until
the source evidence
supports a
specific assignment.

A future taxonomy
may benefit from
greater emphasis
on interaction
mechanism,
control architecture,
and physical
constraints rather
than hardware
appearance alone.

\paragraph*{Limitations of the extracted variables.}

The extraction table supports
the descriptive and
critical synthesis
presented here,
but it cannot
answer every
question relevant
to active-matter
physics.

In particular,
the extracted
agreement labels
do not establish
whether a
theoretical parameter
was measured
independently,
fitted to
experimental data,
or adjusted
after observing
the collective
response.

Similarly,
the risk-of-bias
fields summarize
study-level
limitations but
do not replace
a claim-specific
assessment of
all reported
observables.

The maximum
number of
hardware agents
does not capture
every relevant
measure of
experimental scale.
A microscopic
population may
contain many
particles without
resolving each
trajectory.
A mechanically
connected active
solid may have
many actuators
but only a
limited number
of independent
collective modes.

The theoretical
categories are
also multi-label.
A study comparing
an agent-based
simulation with
an active Brownian
model can appear
in both groups.
Thus, totals
across frameworks
cannot be treated
as independent
study counts.

These limitations
motivate the
primarily narrative,
mechanism-oriented
synthesis used
in this review.

\subsection{What the living design can and cannot fix}
\label{sec:discussion-living}

The living design is intended to keep the evidence
base current as new robotic active-matter
experiments appear. Its principal advantage
is that searches, screening decisions,
extraction records, and synthesis outputs
can be updated under an explicit protocol.

However, living updating does not
automatically improve the reliability
of the underlying literature.
It can correct review-level
errors and incorporate new
evidence. It cannot create
missing experimental
measurements or
independent replications.

\paragraph*{What updating can correct.}

The monthly search schedule can identify
newly indexed experimental studies,
new versions of preprints,
and publications that
were previously inaccessible.

When full texts become available,
abstract-only inclusions
can be reassessed.
This may improve
confidence in
eligibility decisions
and support
more complete
extraction of
experimental methods.

The living workflow
can also reconcile
duplicate records
and publication versions.

The correction of
included-study records
is an example.
Such corrections
change descriptive
counts, even when
the underlying
scientific literature
has not changed.
This demonstrates
why snapshot
versioning matters.

This demonstrates
why snapshot
versioning matters.
A numerical
claim must refer
to a specific
evidence-base
version, and
revisions should
remain traceable.

Future updates
can also resolve
missing platform
assignments,
improve venue
classification,
and refine
the distinction
between preprints
and peer-reviewed
publications.

Where additional
reports provide
new experimental
detail, risk-of-bias
ratings may change.

For example,
an accessible
supplement could
establish independent
replication counts
or provide
previously unavailable
data and code.

Such corrections
can improve
the accuracy
of the review
without changing
the original
physical experiment.

\paragraph*{Updating methodological validation.}

The review's automated
workflow may also
change over time.
Screening prompts,
models, adjudication
rules, and extraction
procedures can
be revised.

These changes
should not be
treated as
administrative
details.
They may alter
which studies
are included
and how evidence
is classified.

A living review
therefore requires
versioned records
of the methods
used at each
snapshot.

The October 17, 2026
validation describes
the workflow evaluated
at that time.
It does not
establish the
performance of
future models
or modified
screening procedures.

Major changes
should trigger
new validation
against independently
assessed records.

The distinction
between blind
and final
reference labels
must remain
explicit in
future reports.
Otherwise, apparent
improvements in
agreement could
reflect changes
in the reference
labels rather
than genuine
changes in
screening performance.

Similarly,
the adjudication
role of the
Opus judge
should remain
separate from
its evaluation
as a blind
screener.

This separation
is essential
for interpreting
the resulting
performance
statistics.

\paragraph*{Updating certainty assessments.}

A living review
should not simply
increase study
counts and
repeat the
same conclusions.

New evidence
should be
examined for
its effect
on the certainty
of specific
scientific claims.

For example,
a new experimental
demonstration of
flocking may add
breadth to the
phenomenon inventory.
Yet, it may
have little effect
on certainty
about a
universal
ordering transition
unless it
provides relevant
quantitative
measurements.

Conversely,
one well-designed
experiment measuring
finite-size behavior
across independent
conditions could
materially change
the assessment
of a specific
theoretical claim.

Likewise,
an independent
measurement of
nonreciprocal
elastic response
could strengthen
the evidence
for active-solid
theory more
than several
additional
demonstrations of
programmable
mechanical motion.

Hence, the
living synthesis
should update
certainty in
response to
the content
of new evidence,
not merely
its volume.

A useful
update trigger
would be
the appearance
of a study
that changes
the balance
of quantitative
agreement,
risk of bias,
or independent
cross-platform
support for
a principal
conclusion.

The same applies
when a new
result contradicts
an established
interpretation.
A living review
must be able
to downgrade
certainty as
well as
upgrade it.

\paragraph*{What updating cannot repair.}

The living design
cannot resolve
a lack of
independent
experimental
repetitions when
those repetitions
were never
performed or
reported.

Nor can it
establish
thermodynamic-limit
behavior from
a collection
of small,
geometrically
similar experiments
unless the new
evidence supplies
the necessary
size-dependent
measurements.

Likewise,
a larger
literature does
not eliminate
publication bias.
If unsuccessful
experiments remain
unpublished,
monthly searches
will not retrieve
them.

The living
workflow also
cannot guarantee
that automated
screening and
extraction are
free from
systematic error.
Revalidation
can estimate
and help
identify errors,
but it cannot
prove the
absence of
shared biases
between models
and adjudicators.

Finally,
the living
design cannot
transform
qualitative
model agreement
into independent
predictive validation.
That requires
experiments
designed to
distinguish
between competing
physical explanations.

\paragraph*{A stable protocol with revisable judgments.}

A successful living review
must balance methodological
consistency against the
need to improve
its classifications
and evidence assessments.

The eligibility
criteria should
remain sufficiently
stable to support
comparison between
snapshots.
Changes in
definitions should
be recorded
as protocol
amendments,
not silently
introduced through
new extraction
practice.

At the same
time, a taxonomy
that no longer
describes the
experimental
literature adequately
should not be
preserved solely
for convenience.

The residual
platform category,
for example,
may justify
a documented
revision that
distinguishes
previously
combined
physical mechanisms.

Such a change
should be accompanied
by consistent
reclassification
of earlier
records where
the available
evidence permits it.

Similarly,
a future extraction
could distinguish
fitted from
predicted
theoretical
parameters,
separate
spontaneous
self-assembly
from prescribed
formation,
or record
independent
experimental
groups.

These refinements
would improve
the scientific
meaning of
the synthesis.
However,
they would
also require
careful versioning
to avoid
mistaking
classification
changes for
changes in
the literature.

\paragraph*{Implications for interpretation.}

The principal contribution of
the living design
is therefore
epistemic rather
than merely
administrative.

It allows
the review
to preserve
an auditable
relationship
between evidence,
classification,
certainty,
and scientific
conclusion.

A new study
can change
the evidence
base.
A new full
text can
resolve an
uncertain
record.
A new
validation exercise
can reveal
limitations in
the screening
workflow.
A new
theoretical comparison
can strengthen
or challenge
a physical
interpretation.

Yet, the
central standard
remains unchanged.
The strength
of a conclusion
depends on
the observations
and tests
that support it,
not on
the frequency
with which
the review
is updated.

\subsection{Overall interpretation}
\label{sec:discussion-overall}

Taken together, the evidence establishes
robotic active matter as a versatile
experimental approach to collective
nonequilibrium physics.

The strongest conclusion is that
physical robotic collectives
can realize a broad range
of organized behaviors
through programmed,
mechanical, and
field-mediated
interactions.

A second conclusion
is that programmability
provides a valuable
experimental means
of changing the
microscopic rules
thought to determine
collective behavior.

However, the theoretical
interpretation of these
experiments requires
greater care than
the visual similarity
of collective states
might suggest.

The evidence is
less decisive
for universal
phase behavior,
large-system scaling,
and the independent
predictive validity
of particular
theoretical models.

The risk-of-bias
assessment identifies
three particularly
important constraints:
incomplete replication
reporting,
limited finite-size
characterization,
and insufficiently
documented quantitative
theory tests.

Incomplete data
and code availability
further restrict
independent verification.

These limitations
are not specific
to one platform.
They occur
throughout the
experimental
taxonomy,
although their
importance differs
between scientific
questions.

The living design
can improve
coverage,
resolve incomplete
records,
correct classifications,
and update
certainty assessments.
It cannot
replace the
experimental
measurements
needed to
distinguish
physical mechanisms.

We therefore
regard the
next stage
of robotic
active-matter
research as
a transition
from realizing
collective
behavior to
testing it
predictively.

The most
informative
experiments
will characterize
microscopic
interactions,
measure collective
observables
with independent
replication,
vary relevant
parameters
systematically,
and test
theoretical
predictions
outside the
conditions used
for fitting.

Cross-platform
replication
will then
help identify
which principles
are genuinely
general and
which depend
on particular
actuation,
sensing,
or mechanical
implementations.

In this
sense,
the value
of robotic
active matter
lies not
only in
constructing
controllable
nonequilibrium
systems.
It lies
in using
that control
to make
physical theories
vulnerable
to precise
experimental
tests.

